\documentclass[11pt,a4paper]{article}

% =============================================================================
% 1. PACKAGES (DO NOT EDIT)
% =============================================================================
\usepackage[utf8]{inputenc}
\usepackage[T1]{fontenc}
\usepackage{lmodern}
\usepackage[a4paper, top=3.5cm, bottom=2.5cm, left=2.5cm, right=2.5cm, headheight=42pt, headsep=1.0cm]{geometry}
\usepackage[table]{xcolor}
\usepackage{tabularx}
\usepackage{array}
\usepackage{fancyhdr}
\usepackage{graphicx}
\usepackage{amsmath}
\usepackage[hidelinks]{hyperref}

% =============================================================================
% 2. DOCUMENT METADATA —>>> PUT PROJECT DATA HERE <<<—
% (Values entered here will automatically populate Page 1 and all page headers)
% =============================================================================
\newcommand{\Organization}{HSB}                     % <-- Put organization name here (e.g., HSB)
\newcommand{\ProjectName}{EWAS} % <-- Put project name here (remove \{ \})
\newcommand{\DocTitle}{Space Segment - Technical Requirements Specification}  % <-- Put document title here
\newcommand{\DocNumber}{EWAS-HSB-SSEG-TRS-03 0.1}         % <-- Put document number here (DOC-NO)
\newcommand{\DocIssue}{0.1}           % <-- Put issue/version number here (e.g., 1.0)
\newcommand{\DocDate}{04.11.2026}             % <-- Put issue date here (e.g., 2026-11-04)
\newcommand{\CINumber}{03}       % <-- Put Configuration Item (CI) number here

% =============================================================================
% 3. COLORS AND TEMPLATE STYLING (DO NOT EDIT)
% =============================================================================
\definecolor{TitleBlue}{HTML}{1F497D}         % Main title color on Page 1
\definecolor{HeadingBlue}{HTML}{365F91}       % Section heading color (H1)
\definecolor{SubHeadingBlue}{HTML}{4F81BD}    % Subsection heading (H2) and rule color
\definecolor{TableHeaderBlue}{HTML}{4F81BD}   % Table header background color
\definecolor{TableRowLightBlue}{HTML}{DBE5F1} % Alternating table row background color
\definecolor{TableBorderBlue}{HTML}{95B3D7}   % Styled table border color

% Paragraph settings
\setlength{\parindent}{0pt}
\setlength{\parskip}{0.6em}

% Disable section numbering (matching the PDF template) while keeping TOC entries.
% To enable section numbering (e.g., 1, 1.1), change -1 to 3.
\setcounter{secnumdepth}{-1}
\setcounter{tocdepth}{2}

% Section, subsection, and list title styling
\makeatletter
\renewcommand\section{\@startsection{section}{1}{\z@}%
    {-3.0ex \@plus -1ex \@minus -.2ex}%
    {1.5ex \@plus .2ex}%
    {\normalfont\Large\bfseries\sffamily\color{HeadingBlue}}}

\renewcommand\subsection{\@startsection{subsection}{2}{\z@}%
    {-2.5ex \@plus -1ex \@minus -.2ex}%
    {1.0ex \@plus .2ex}%
    {\normalfont\large\bfseries\sffamily\color{SubHeadingBlue}}}

\renewcommand\subsubsection{\@startsection{subsubsection}{3}{\z@}%
    {-2.0ex \@plus -1ex \@minus -.2ex}%
    {0.8ex \@plus .2ex}%
    {\normalfont\normalsize\bfseries\sffamily\color{SubHeadingBlue}}}

% Dotted leaders in the Table of Contents for top-level sections as well
\renewcommand*\l@section{\@dottedtocline{1}{0em}{1.5em}}
\renewcommand*\l@subsection{\@dottedtocline{2}{1.5em}{2.3em}}

% Centered blue headings for Table of Contents, List of Tables, and List of Figures
\renewcommand{\contentsname}{Table of Contents}
\renewcommand{\listtablename}{List of Tables}
\renewcommand{\listfigurename}{List of Figures}

\renewcommand\tableofcontents{
    {\centering\normalfont\Large\bfseries\sffamily\color{HeadingBlue}\contentsname\par}
    \vspace{0.8em}
    \@starttoc{toc}
}
\renewcommand\listoftables{
    {\centering\normalfont\Large\bfseries\sffamily\color{HeadingBlue}\listtablename\par}
    \vspace{0.8em}
    \@starttoc{lot}
}
\renewcommand\listoffigures{
    {\centering\normalfont\Large\bfseries\sffamily\color{HeadingBlue}\listfigurename\par}
    \vspace{0.8em}
    \@starttoc{lof}
}
\makeatother

% Helper macros
\newcommand{\thead}[1]{\textcolor{white}{\textbf{\textsf{#1}}}}
\providecommand{\abbr}[1]{#1}

% =============================================================================
% 4. HEADER AND FOOTER (WITH UNIVERSITY LOGOS)
% =============================================================================
\pagestyle{fancy}
\fancyhf{}
\renewcommand{\headrulewidth}{0pt}
\renewcommand{\footrulewidth}{0pt}

\fancyhead[L]{
    \raisebox{-0.5\height}{\includegraphics[height=1.05cm]{figs/hsb_logo.png}}
}
\fancyhead[C]{
    {\small\textbf{Early Warning Alarm System (\abbr{EWAS})}}
}
\fancyhead[R]{
    \raisebox{-0.5\height}{\includegraphics[height=1.2cm]{figs/pg_logo.jpg}}
}
\fancyfoot[R]{\thepage}

% Style for Page 1 (logos in top corners, no page number in footer)
\fancypagestyle{firstpage}{
    \fancyhf{}
    \renewcommand{\headrulewidth}{0pt}
    \renewcommand{\footrulewidth}{0pt}
    \fancyhead[L]{\raisebox{-0.5\height}{\includegraphics[height=1.05cm]{figs/hsb_logo.png}}}
    \fancyhead[R]{\raisebox{-0.5\height}{\includegraphics[height=1.2cm]{figs/pg_logo.jpg}}}
}

% =============================================================================
% DOCUMENT BODY
% =============================================================================
\begin{document}

% -----------------------------------------------------------------------------
% PAGE 1: TITLE PAGE & DOCUMENT INFORMATION
% -----------------------------------------------------------------------------
\thispagestyle{firstpage} % Use \thispagestyle{empty} if you do not want logos on Page 1
\phantomsection
\addcontentsline{toc}{section}{Document Information}

\begin{center}
    {\large \ProjectName}
\end{center}

\vspace*{5.5cm}

\begin{center}
    {\fontsize{24}{28}\selectfont\sffamily\color{TitleBlue}\DocTitle\par}
    \vspace{0.35cm}
    {\color{SubHeadingBlue}\hrule height 0.8pt}
\end{center}

\vfill

% --- HOW TO EDIT THE PAGE 1 TABLE ---
% Part 1 (first 4 rows): automatically populated from Section 2 at the top of the file.
% Part 2 (signatures & roles): enter data between the & symbols in the format:
%     {Full Name} & {Role Name} & {Signature/Date (or leave blank)} \\[1.2em] \hline
{\renewcommand{\arraystretch}{1.25}%
\arrayrulecolor{black}%
\noindent
\begin{tabularx}{\textwidth}{|X|X|}
    \hline
    Document Number & \DocNumber \\ \hline
    Issue           & \DocIssue  \\ \hline
    Issue Date      & \DocDate   \\ \hline
    CI Number       & \CINumber  \\
\end{tabularx}\par\nointerlineskip
\noindent
\begin{tabularx}{\textwidth}{|X|X|X|}
    \hline
    % COLUMN 1: Name             & COLUMN 2: Responsibility & COLUMN 3: Signature/Date
    Name                         & Responsibility           & Signature/Date                   \\ \hline
    Prepared by                  &                          &                                  \\ \hline
    {Jan Krajewski}              & {Author}     & {Krajewski 04.11.2026} \\[1.2em] \hline
    Checked by                   &                          &                                  \\ \hline
    {Sergiusz Pyskowacki}              & {Reviewer}     & {Pyskowacki 04.11.2026} \\[1.2em] \hline
    Approved by                  &                          &                                  \\ \hline
    {Put name here}              & {Approver}     & {Put Date/Signature or leave empty} \\[1.2em] \hline
    {Put Name here}              & {Put Role Name here}     & {Put Date/Signature or leave empty} \\[1.2em] \hline
\end{tabularx}}

\newpage

% -----------------------------------------------------------------------------
% PAGE 2: DOCUMENT CHANGE RECORD
% -----------------------------------------------------------------------------
\section{Document Change Record}

% --- HOW TO EDIT THE DOCUMENT CHANGE RECORD TABLE ---
% Each row has 4 columns separated by &:
% \textbf{Issue} & {Date} & {Description} & {Sections Affected} \\ \hline
% To add a new row, copy the last line before \end{tabularx}.
{\renewcommand{\arraystretch}{1.25}%
\arrayrulecolor{TableBorderBlue}%
\rowcolors{2}{white}{TableRowLightBlue}%
\noindent
\begin{tabularx}{\textwidth}{|p{2.0cm}|p{2.5cm}|X|p{4.5cm}|}
    \hline
    \rowcolor{TableHeaderBlue}
    % COLUMN 1: Issue & COLUMN 2: Date  & COLUMN 3: Description      & COLUMN 4: Sections Affected
    \thead{Issue}     & \thead{Date}    & \thead{Description}        & \thead{Sections Affected}     \\ \hline
    \textbf{0.1}      & {Put Date here} & Initial Issue              & Initial Draft                 \\ \hline
    \textbf{1.0}      & {Put Date here} & First Release              & All                           \\ \hline
    % \textbf{1.1}    & {Put Date here} & {Put change description}   & {Put affected sections}       \\ \hline
\end{tabularx}}

\newpage

% -----------------------------------------------------------------------------
% PAGE 3: TABLE OF CONTENTS, LIST OF TABLES, LIST OF FIGURES (AUTO-GENERATED)
% -----------------------------------------------------------------------------
\tableofcontents

\vspace{2.5cm}
\listoftables

\vspace{2.5cm}
\listoffigures

\newpage

% -----------------------------------------------------------------------------
% PAGE 4: REFERENCES (AD, RD, TERMS & ABBREVIATIONS)
% -----------------------------------------------------------------------------
\section{References}

\subsection{Applicable Documents (AD)}

% --- HOW TO EDIT THE APPLICABLE DOCUMENTS (AD) TABLE ---
% Columns are separated by &:
% \textbf{AD-XX} & {Document Number} & {Document Title} & {Issue} \\ \hline
{\renewcommand{\arraystretch}{1.25}%
\arrayrulecolor{TableBorderBlue}%
\rowcolors{2}{white}{TableRowLightBlue}%
\noindent
\begin{tabularx}{\textwidth}{|p{2.4cm}|p{3.6cm}|X|p{2.4cm}|}
    \hline
    \rowcolor{TableHeaderBlue}
    % COLUMN 1: AD No. & COLUMN 2: Document Number & COLUMN 3: Title         & COLUMN 4: Issue
    \thead{AD No.}     & \thead{Document Number}   & \thead{Title}           & \thead{Issue}     \\ \hline
    \textbf{AD-01}     & {Put Doc. Number here}    & {Put Title here}        & {Put Issue here}  \\ \hline
    \textbf{AD-02}     & {Put Doc. Number here}    & {Put Title here}        & {Put Issue here}  \\ \hline
    \textbf{AD-03}     & {Put Doc. Number here}    & {Put Title here}        & {Put Issue here}  \\ \hline
    \textbf{AD-04}     & {Put Doc. Number here}    & {Put Title here}        & {Put Issue here}  \\ \hline
    \textbf{AD-05}     & {Put Doc. Number here}    & {Put Title here}        & {Put Issue here}  \\ \hline
\end{tabularx}}

\vspace{0.8em}
\subsection{Reference Documents (RD)}

% --- HOW TO EDIT THE REFERENCE DOCUMENTS (RD) TABLE ---
% Columns are separated by &:
% \textbf{RD-XX} & {Document Number} & {Document Title} & {Issue} \\ \hline
{\renewcommand{\arraystretch}{1.25}%
\arrayrulecolor{TableBorderBlue}%
\rowcolors{2}{white}{TableRowLightBlue}%
\noindent
\begin{tabularx}{\textwidth}{|p{2.4cm}|p{3.6cm}|X|p{2.4cm}|}
    \hline
    \rowcolor{TableHeaderBlue}
    % COLUMN 1: RD No. & COLUMN 2: Document Number & COLUMN 3: Title         & COLUMN 4: Issue
    \thead{RD No.}     & \thead{Document Number}   & \thead{Title}           & \thead{Issue}     \\ \hline
    \textbf{RD-01}     & {Put Doc. Number here}    & {Put Title here}        & {Put Issue here}  \\ \hline
    \textbf{RD-02}     & {Put Doc. Number here}    & {Put Title here}        & {Put Issue here}  \\ \hline
    \textbf{RD-03}     & {Put Doc. Number here}    & {Put Title here}        & {Put Issue here}  \\ \hline
    \textbf{RD-04}     & {Put Doc. Number here}    & {Put Title here}        & {Put Issue here}  \\ \hline
    \textbf{RD-05}     & {Put Doc. Number here}    & {Put Title here}        & {Put Issue here}  \\ \hline
\end{tabularx}}

\vspace{0.8em}
\subsection{Terms, Definitions and Abbreviations}

% --- HOW TO EDIT THE ABBREVIATIONS TABLE ---
% Columns are separated by &:
% \textbf{ABBREVIATION} & {Full Meaning / Definition} \\ \hline
{\renewcommand{\arraystretch}{1.25}%
\arrayrulecolor{TableBorderBlue}%
\rowcolors{2}{white}{TableRowLightBlue}%
\noindent
\begin{tabularx}{\textwidth}{|X|X|}
    \hline
    \rowcolor{TableHeaderBlue}
    % COLUMN 1: Abbreviation          & COLUMN 2: Meaning
    \thead{Abbreviation}              & \thead{Meaning}                                   \\ \hline
    \textbf{ECSS}                     & European Cooperation for Space Standardization    \\ \hline
    \textbf{CI}                       & Configuration Item                                \\ \hline
    \textbf{{Put Abbreviation here}}  & {Put Meaning / Definition here}                   \\ \hline
\end{tabularx}}

\newpage

% -----------------------------------------------------------------------------
% PAGE 5: SCOPE
% -----------------------------------------------------------------------------
\section{Scope}

% ---> PUT DATA HERE (SCOPE SECTION CONTENT) <---
{Put Scope content here.}

\newpage

% -----------------------------------------------------------------------------
% PAGE 6: INTRODUCTION
% -----------------------------------------------------------------------------
\section{Introduction}

% ---> PUT DATA HERE (INTRODUCTION SECTION CONTENT) <---
{Put Introduction content here.}

\newpage

% -----------------------------------------------------------------------------
% PAGE 7: MAIN CONTENT
% -----------------------------------------------------------------------------
\section{Main Content}

% ---> PUT DATA HERE (MAIN CONTENT OF THE DOCUMENT) <---
{Put Main Content here.}

% EXAMPLE OF ADDING A SUBSECTION (uncomment if needed):
% \subsection{Put Subsection Title here}
% {Put subsection content here.}

\newpage

% -----------------------------------------------------------------------------
% PAGE 8: ANNEX A
% -----------------------------------------------------------------------------
\section{Annex A}

% ---> PUT DATA HERE (ANNEX A CONTENT) <---
{Put Annex A content here.}

\newpage

% -----------------------------------------------------------------------------
% PAGE 9: ANNEX B
% -----------------------------------------------------------------------------
\section{Annex B}

% ---> PUT DATA HERE (ANNEX B CONTENT) <---
{Put Annex B content here.}

\end{document}
